\chapter{Learning Theory — Hedge}
%%% generated by the blueprint pipeline from TCSlib.LearningTheory.Hedge
\section{Overview}

This module formalizes the Hedge (Multiplicative Weights) algorithm in the expert
setting with $N$ experts and $T$ rounds.  It establishes the standard regret bound
$(\ln N)/\eta + \eta T/2$ (and the tighter $(\ln N)/\eta + \eta T/8$ variant via
Hoeffding's lemma), together with the no-regret corollary that average regret
vanishes as $T \to \infty$ when $\eta = \sqrt{2 \ln N / T}$.

%%%%%%%%%%%%%%%%%%%%%%%%%%%%%%%%%%%%%%%%%%%%%%%%%%%%%%%%%%%%%%%%%%%%%%%%%%%%%%%
\section{Declarations}

\begin{definition}[Loss sequence]
\lean{LossSeq}
\label{LossSeq}
\leanok
A \emph{loss sequence} $\ell$ for $N$ experts over $T$ rounds is a function
$\ell : \mathrm{Fin}\,T \to \mathrm{Fin}\,N \to \mathbb{R}$, i.e., at each round
$t < T$ the adversary reveals a loss vector $\ell_t : \mathrm{Fin}\,N \to \mathbb{R}$
whose values are intended to lie in $[0,1]$.
\end{definition}

\begin{definition}[Valid loss sequence]
\lean{LossSeq.Valid}
\label{LossSeq.Valid}
\leanok
\uses{LossSeq}
A loss sequence $\ell$ is \emph{valid} if all individual losses are in $[0,1]$:
$\forall t\,i,\; 0 \le \ell_t(i) \le 1$.
\end{definition}

\begin{definition}[Cumulative loss]
\lean{cumLoss}
\label{cumLoss}
\leanok
\uses{LossSeq}
The \emph{cumulative loss} of expert $i$ through the first $t$ rounds is
\[
  L_t(i) \;=\; \sum_{s < t} \ell_s(i),
\]
where the sum ranges over all rounds $s \in \mathrm{Fin}\,T$ with $s < t$.
\end{definition}

\begin{definition}[Hedge weight]
\lean{hedgeWeight}
\label{hedgeWeight}
\leanok
\uses{LossSeq, cumLoss}
The \emph{unnormalized Hedge weight} of expert $i$ at time $t$ with learning
rate $\eta$ is
\[
  w_t(i) \;=\; \exp\!\bigl(-\eta \cdot L_t(i)\bigr).
\]
Experts with smaller cumulative loss receive larger weight.
\end{definition}

\begin{definition}[Potential]
\lean{potential}
\label{potential}
\leanok
\uses{LossSeq, hedgeWeight}
The \emph{potential} (sum of unnormalized weights) at time $t$ is
\[
  W_t \;=\; \sum_{i=1}^{N} w_t(i).
\]
The potential telescopes to relate the learner's cumulative loss to the best
expert's loss.
\end{definition}

\begin{definition}[Hedge distribution]
\lean{hedgeDist}
\label{hedgeDist}
\leanok
\uses{LossSeq, hedgeWeight, potential}
The \emph{Hedge distribution} at round $t$ is the normalization of the weight
vector:
\[
  p_t(i) \;=\; \frac{w_t(i)}{W_t}.
\]
\end{definition}

\begin{definition}[Expected loss of Hedge]
\lean{hedgeLoss}
\label{hedgeLoss}
\leanok
\uses{LossSeq, hedgeDist}
The \emph{expected loss} of the learner at round $t$ under the Hedge
distribution is
\[
  \widehat{\ell}_t \;=\; \sum_{i=1}^{N} p_t(i)\,\ell_t(i).
\]
\end{definition}

\begin{definition}[Cumulative loss of Hedge]
\lean{hedgeCumLoss}
\label{hedgeCumLoss}
\leanok
\uses{LossSeq, hedgeLoss}
The \emph{cumulative loss} of the Hedge algorithm over all $T$ rounds is
\[
  \widehat{L}_T \;=\; \sum_{t=0}^{T-1} \widehat{\ell}_t.
\]
\end{definition}

\begin{definition}[Best expert loss]
\lean{bestExpertLoss}
\label{bestExpertLoss}
\leanok
\uses{LossSeq, cumLoss}
The \emph{best expert loss in hindsight} is the minimum cumulative loss
achieved by any single expert:
\[
  L^* \;=\; \min_{i \in [N]} L_T(i) \;=\; \inf_{i : \mathrm{Fin}\,N} \,L_T(i).
\]
\end{definition}

\begin{definition}[Regret]
\lean{regret}
\label{regret}
\leanok
\uses{LossSeq, bestExpertLoss, hedgeCumLoss}
The \emph{regret} of the Hedge algorithm is the excess cumulative loss over the
best expert in hindsight:
\[
  R_T(\eta) \;=\; \widehat{L}_T - L^*.
\]
\end{definition}

\begin{lemma}[Initial value of the potential]
\lean{potential_zero}
\label{potential_zero}
\leanok
\uses{LossSeq, cumLoss, hedgeWeight, potential}
\difficulty{3}
Fix $N \ge 1$ experts, a horizon $T$, and a loss sequence $\ell$ over these experts and
rounds, and let $\eta \in \bbr$ be a learning rate. At time $0$ every expert has zero
cumulative loss, so each unnormalized Hedge weight equals $1$, and the potential is
\[
  W_0 \;=\; \sum_{i=1}^{N} \exp\!\bigl(-\eta \cdot L_0(i)\bigr) \;=\; N .
\]
\end{lemma}

\begin{lemma}[Positivity of Hedge weights]
\lean{hedgeWeight_pos}
\label{hedgeWeight_pos}
\leanok
\uses{LossSeq, hedgeWeight}
\difficulty{1}
Fix a learning rate $\eta \in \bbr$, a loss sequence $\ell$ for $N$ experts over $T$
rounds, a time $t$, and an expert $i$. Then the unnormalized Hedge weight $w_t(i) =
\exp\!\bigl(-\eta \cdot L_t(i)\bigr)$, where $L_t(i) = \sum_{s < t} \ell_s(i)$ is the
cumulative loss of expert $i$ through the first $t$ rounds, is strictly positive.
\end{lemma}

\begin{lemma}[Positivity of the Hedge potential]
\lean{potential_pos}
\label{potential_pos}
\leanok
\uses{LossSeq, hedgeWeight_pos, potential}
\difficulty{2}
Fix a learning rate $\eta \in \bbr$, a loss sequence $\ell$ for $N \ge 1$ experts over
$T$ rounds, and a time $t$. Then the potential at time $t$ is strictly positive:
\[
  W_t \;=\; \sum_{i=1}^{N} \exp\!\bigl(-\eta \cdot L_t(i)\bigr) \;>\; 0,
\]
where $L_t(i) = \sum_{s < t} \ell_s(i)$ is the cumulative loss of expert $i$ through the
first $t$ rounds.
\end{lemma}

\begin{lemma}[Chord bound for the exponential on the unit interval]
\lean{exp_neg_le_linear}
\label{exp_neg_le_linear}
\leanok
\difficulty{4}
For every $\eta > 0$ and every $x \in [0,1]$, the value $e^{-\eta x}$ lies below the
chord joining the endpoints of the graph of $t \mapsto e^{-\eta t}$ over $[0,1]$; that
is, the inequality
\[
  e^{-\eta x} \;\le\; 1 - \bigl(1 - e^{-\eta}\bigr)\,x
\]
holds, the right-hand side being the affine function equal to $1$ at $x = 0$ and to
$e^{-\eta}$ at $x = 1$.
\end{lemma}

\begin{lemma}[Cumulative loss recurrence]
\lean{cumLoss_succ}
\label{cumLoss_succ}
\leanok
\uses{LossSeq, cumLoss}
\difficulty{4}
Let $\ell$ be a loss sequence for $N$ experts over $T$ rounds, and write $L_t(i) =
\sum_{s < t} \ell_s(i)$ for the cumulative loss of expert $i$ through the first $t$
rounds. Then for every round $t < T$ and every expert $i$,
\[
  L_{t+1}(i) \;=\; L_t(i) + \ell_t(i).
\]
\end{lemma}

\begin{lemma}[Cumulative loss at the horizon]
\lean{cumLoss_horizon}
\label{cumLoss_horizon}
\leanok
\uses{LossSeq, cumLoss}
\difficulty{2}
Fix $N$ experts and a horizon $T$, and let $\ell$ be a loss sequence assigning to each
round $t$ and expert $i$ a real loss $\ell_t(i)$. For every expert $i$, the cumulative
loss through the first $T$ rounds equals the total loss over all rounds: $L_T(i) =
\sum_{t} \ell_t(i)$, where $t$ ranges over all $T$ rounds.
\end{lemma}

\begin{lemma}[Recursion for the Hedge weight]
\lean{hedgeWeight_succ}
\label{hedgeWeight_succ}
\leanok
\uses{LossSeq, cumLoss_succ, hedgeWeight}
\difficulty{2}
Fix a learning rate $\eta \in \bbr$ and a loss sequence $\ell$ for $N$ experts over $T$
rounds. For every round $t \in \mathrm{Fin}\,T$ and every expert $i \in
\mathrm{Fin}\,N$, the unnormalized Hedge weights satisfy
\[
  w_{t+1}(i) \;=\; w_t(i)\cdot \exp\!\bigl(-\eta\,\ell_t(i)\bigr).
\]
\end{lemma}

\begin{lemma}[The Hedge distribution sums to one]
\lean{hedgeDist_sum}
\label{hedgeDist_sum}
\leanok
\uses{LossSeq, hedgeDist, potential_pos}
\difficulty{2}
Fix a number of experts $N \ge 1$, a horizon $T$, a learning rate $\eta \in \bbr$, a
loss sequence $\ell$ for $N$ experts over $T$ rounds, and a round $t$. Then the Hedge
distribution at round $t$ sums to one over all experts,
\[
  \sum_{i=1}^{N} p_t(i) = 1,
\]
so it is a probability distribution on the $N$ experts.
\end{lemma}

\begin{lemma}[Potential ratio bound for Hedge]
\lean{potential_ratio_le}
\label{potential_ratio_le}
\leanok
\uses{LossSeq, LossSeq.Valid, exp_neg_le_linear, hedgeDist, hedgeLoss, hedgeWeight, hedgeWeight_pos, hedgeWeight_succ, potential, potential_pos}
\difficulty{5}
Fix $N \ge 1$ experts, a horizon of $T$ rounds, a learning rate $\eta > 0$, and a valid
loss sequence $\ell$. Then for every round $t$ with $t < T$, the consecutive potentials
satisfy
\[
  \frac{W_{t+1}}{W_t} \;\le\; 1 - \bigl(1 - e^{-\eta}\bigr)\,\widehat{\ell}_t,
\]
where $\widehat{\ell}_t = \sum_{i=1}^{N} p_t(i)\,\ell_t(i)$ is the expected loss of the
learner under the Hedge distribution at round $t$.
\end{lemma}

\begin{lemma}[One-step logarithmic potential bound for Hedge]
\lean{log_potential_step}
\label{log_potential_step}
\leanok
\uses{LossSeq, LossSeq.Valid, hedgeLoss, potential, potential_pos, potential_ratio_le}
\difficulty{4}
Fix $N \ge 1$ experts and $T$ rounds, a learning rate $\eta > 0$, and a valid loss
sequence $\ell$, so that $0 \le \ell_t(i) \le 1$ for every round $t$ and expert $i$.
Writing $W_t = \sum_{i=1}^N \exp\!\bigl(-\eta L_t(i)\bigr)$ for the Hedge potential and
$\widehat{\ell}_t = \sum_{i=1}^N p_t(i)\,\ell_t(i)$ for the expected loss under the
Hedge distribution, the one-round change in log-potential is controlled by the expected
loss: for each round $t$,
\[
  \ln W_{t+1} - \ln W_t \;\le\; -\bigl(1 - e^{-\eta}\bigr)\,\widehat{\ell}_t.
\]
\end{lemma}

\begin{lemma}[Quadratic upper bound on $e^{-t}$]
\lean{exp_neg_le_quadratic}
\label{exp_neg_le_quadratic}
\leanok
\difficulty{4}
For every real number $t \ge 0$, the value $e^{-t}$ is at most the second-order Taylor
polynomial $1 - t + t^2/2$; that is,
\[
  e^{-t} \;\le\; 1 - t + \frac{t^2}{2}.
\]
\end{lemma}

\begin{lemma}[Lower bound on $1 - e^{-\eta}$]
\lean{one_sub_exp_neg_ge}
\label{one_sub_exp_neg_ge}
\leanok
\uses{exp_neg_le_quadratic}
\difficulty{2}
For every real number $\eta > 0$, the quantity $1 - e^{-\eta}$ is bounded below by its
second-order Taylor expansion at the origin:
\[
  1 - e^{-\eta} \;\ge\; \eta - \frac{\eta^{2}}{2}.
\]
\end{lemma}

\begin{lemma}[Potential lower bound via the best expert]
\lean{potential_ge_best_expert}
\label{potential_ge_best_expert}
\leanok
\uses{LossSeq, bestExpertLoss, cumLoss, hedgeWeight, potential}
\difficulty{4}
Fix $N \ge 1$ experts and $T$ rounds, a loss sequence $\ell$, and a learning rate $\eta
> 0$. Then the final potential is at least the unnormalized Hedge weight of the best
expert in hindsight:
\[
  W_T \;\ge\; \exp\!\bigl(-\eta \, L^*\bigr),
\]
where $L^* = \min_{i} L_T(i)$ is the smallest cumulative loss over the full horizon $T$.
\end{lemma}

\begin{lemma}[Nonnegativity of the Hedge distribution]
\lean{hedgeDist_nonneg}
\label{hedgeDist_nonneg}
\leanok
\uses{LossSeq, hedgeDist, hedgeWeight_pos, potential_pos}
\difficulty{1}
Fix $N \ge 1$ experts and $T$ rounds, a learning rate $\eta \in \bbr$, and a loss
sequence $\ell$. For every time $t$ and every expert $i$, the Hedge distribution is
nonnegative, that is, $p_t(i) \ge 0$.
\end{lemma}

\begin{lemma}[Hedge's expected loss is at most one]
\lean{hedgeLoss_le_one}
\label{hedgeLoss_le_one}
\leanok
\uses{LossSeq, LossSeq.Valid, hedgeDist, hedgeDist_nonneg, hedgeDist_sum, hedgeLoss}
\difficulty{3}
Fix a learning rate $\eta \in \bbr$, and let $\ell$ be a valid loss sequence for $N \ge
1$ experts over $T$ rounds, so that $0 \le \ell_t(i) \le 1$ for every round $t$ and
expert $i$. Then at each round $t$ the expected loss of the learner under the Hedge
distribution satisfies $\widehat{\ell}_t \le 1$.
\end{lemma}

\begin{lemma}[Nonnegativity of Hedge's expected loss]
\lean{hedgeLoss_nonneg}
\label{hedgeLoss_nonneg}
\leanok
\uses{LossSeq, LossSeq.Valid, hedgeDist_nonneg, hedgeLoss}
\difficulty{2}
Let $\ell$ be a valid loss sequence for $N \ge 1$ experts over $T$ rounds, fix any
learning rate $\eta \in \bbr$, and let $t$ be one of the rounds. Then the learner's
expected loss $\widehat{\ell}_t = \sum_{i=1}^{N} p_t(i)\,\ell_t(i)$ under the Hedge
distribution is nonnegative, $\widehat{\ell}_t \ge 0$.
\end{lemma}

\begin{theorem}[Regret bound for the Hedge algorithm]
\lean{hedge_regret_bound}
\label{hedge_regret_bound}
\leanok
\uses{LossSeq, LossSeq.Valid, bestExpertLoss, hedgeCumLoss, hedgeLoss, hedgeLoss_le_one, hedgeLoss_nonneg, log_potential_step, one_sub_exp_neg_ge, potential, potential_ge_best_expert, potential_zero, regret}
\difficulty{5}
Fix $N \ge 1$ experts and $T$ rounds, and let $\ell$ be a valid loss sequence, so that
$0 \le \ell_t(i) \le 1$ for every round $t$ and expert $i$. Then for any learning rate
$\eta$ with $0 < \eta \le 1$, the regret of the Hedge algorithm satisfies
\[
  R_T(\eta) \;\le\; \frac{\ln N}{\eta} + \frac{\eta\,T}{2}.
\]
\end{theorem}

\begin{definition}[Optimal learning rate]
\lean{optimalEta}
\label{optimalEta}
\leanok
The \emph{optimal learning rate} balancing the two terms of the Hedge regret
bound is
\[
  \eta^* \;=\; \sqrt{\frac{2\ln N}{T}}.
\]
\end{definition}

\begin{theorem}[Optimal-rate regret bound for Hedge]
\lean{hedge_regret_optimal}
\label{hedge_regret_optimal}
\leanok
\uses{LossSeq, LossSeq.Valid, hedge_regret_bound, optimalEta, regret}
\difficulty{5}
Consider the Hedge algorithm run on $N > 1$ experts over $T > 0$ rounds against a valid
loss sequence $\ell$, that is, one with $0 \le \ell_t(i) \le 1$ for every round $t$ and
expert $i$. Suppose $2\ln N \le T$, and run Hedge with the learning rate $\eta^* =
\sqrt{2\ln N / T}$. Then the regret—the learner's total expected loss minus the loss of
the best single expert in hindsight—satisfies
\[
  R_T(\eta^*) \;\le\; \sqrt{2\,T\ln N}.
\]
\end{theorem}

\begin{theorem}[Hedge achieves vanishing average regret]
\lean{hedge_no_regret}
\label{hedge_no_regret}
\leanok
\uses{LossSeq, LossSeq.Valid, hedge_regret_optimal, optimalEta, regret}
\difficulty{4}
Fix $N > 1$ experts and $T$ rounds with $T > 0$, and suppose $2\ln N \le T$. Let $\ell$
be a valid loss sequence for $N$ experts over $T$ rounds, so that every individual loss
satisfies $0 \le \ell_t(i) \le 1$. Running the Hedge algorithm with the optimal learning
rate $\eta^* = \sqrt{2\ln N / T}$, the average regret against the best expert in
hindsight is bounded by
\[
  \frac{R_T(\eta^*)}{T} \;\le\; \sqrt{\frac{2\ln N}{T}}.
\]
\end{theorem}

\begin{lemma}[Hoeffding's lemma]
\lean{hoeffding_lemma}
\label{hoeffding_lemma}
\leanok
\uses{bernoulli_mgf_bound}
\difficulty{2}
Let $p \in [0,1]$ and let $h \in \bbr$ be arbitrary. Then
\[
  \log\!\bigl((1-p) + p\,e^{h}\bigr) \;\le\; p\,h + \frac{h^{2}}{8}.
\]
\end{lemma}

\begin{lemma}[Tight per-step logarithmic potential bound for Hedge]
\lean{log_potential_step_tight}
\label{log_potential_step_tight}
\leanok
\uses{LossSeq, LossSeq.Valid, hedgeLoss, hedgeLoss_le_one, hedgeLoss_nonneg, hoeffding_lemma, potential, potential_pos, potential_ratio_le}
\difficulty{4}
Fix $N \ge 1$ experts and a learning rate $\eta > 0$, and let $\ell$ be a valid loss
sequence over $T$ rounds, so that every individual loss satisfies $0 \le \ell_t(i) \le
1$. Write $W_t = \sum_{i=1}^{N} \exp(-\eta L_t(i))$ for the Hedge potential after $t$
rounds, and $\widehat{\ell}_t$ for the expected loss of the learner at round $t$ under
the Hedge distribution. Then for every round $t$ (with $0 \le t < T$),
\[
  \log W_{t+1} - \log W_t \;\le\; -\eta\,\widehat{\ell}_t + \frac{\eta^2}{8}.
\]
\end{lemma}

\begin{theorem}[Tight Hedge regret bound]
\lean{hedge_regret_bound_tight}
\label{hedge_regret_bound_tight}
\leanok
\uses{LossSeq, LossSeq.Valid, bestExpertLoss, hedgeCumLoss, hedgeLoss, log_potential_step_tight, potential, potential_ge_best_expert, potential_zero, regret}
\difficulty{5}
Fix $N \ge 1$ experts and a horizon of $T$ rounds, and run the Hedge algorithm with a
learning rate $\eta > 0$. For every valid loss sequence $\ell$ — one whose losses
satisfy $0 \le \ell_t(i) \le 1$ for all rounds $t$ and experts $i$ — the regret of Hedge
is bounded by
\[
  R_T(\eta) \;\le\; \frac{\log N}{\eta} + \frac{\eta\,T}{8}.
\]
\end{theorem}

\begin{definition}[Tight optimal learning rate]
\lean{optimalEtaTight}
\label{optimalEtaTight}
\leanok
The learning rate that optimizes the tight Hedge bound is
\[
  \eta_{\mathrm{tight}} \;=\; \sqrt{\frac{8\ln N}{T}}.
\]
\end{definition}

\begin{theorem}[Optimal-rate regret bound for Hedge]
\lean{hedge_regret_tight_optimal}
\label{hedge_regret_tight_optimal}
\leanok
\uses{LossSeq, LossSeq.Valid, hedge_regret_bound_tight, optimalEtaTight, regret}
\difficulty{4}
Consider the Hedge algorithm run against a valid loss sequence for $N > 1$ experts over
$T \ge 1$ rounds, so that every individual loss lies in $[0,1]$. If the learning rate is
set to the tight optimal value $\eta_{\mathrm{tight}} = \sqrt{8\ln N / T}$, then the
regret against the best expert in hindsight satisfies
\[
  R_T(\eta_{\mathrm{tight}}) \;\le\; \sqrt{\tfrac{T}{2}\,\ln N}.
\]
\end{theorem}
